\subsection{Genuine compactness of proper native spaces over bounded coefficients}
\label{subsec:proper-native-compactness}

The prescribed seminorm $v$ on the arbitrary commutative ring $R$ is
explicitly bounded on every coefficient. The actual structure morphism
$f:X\to\operatorname{Spec}R$ is proper. Compactness below uses the true
valuative centers and all-coordinate bounded affine spectra. No field,
completeness, Noetherian condition, finite type of a section ring or
Hausdorff conclusion is assumed. This bounded model-base scope is retained;
an arbitrary unbounded prescribed field norm does not supply a coefficient
map into the chosen residue-field valuation ring by this construction.

\begin{lemma}[Actual compact bounded affine representatives]
\label{scheme:compact-bounded-native-affine-images}
\lean{EnvelopeProperty.Berkovich.RelativeSchemeValuedPoint.boundedAffineImage,
EnvelopeProperty.Berkovich.RelativeSchemeValuedPoint.isCompact_boundedAffineImage,
EnvelopeProperty.Berkovich.RelativeSchemeValuedPoint.mem_boundedAffineImage_of_center_mem}
\leanok
\uses{scheme:relative-affine-open-homeomorphism,scheme:proper-center-local-bounds,
spectrum:compact-relative-bounds}
For an actual affine open $U$, take the actual image in the native
prescribed-base space of its relative section-ring spectrum whose every
element value is at most one. That image is compact. Every native point
whose true model center belongs to $U$ lies in this image.
\end{lemma}
\begin{proof}
Every-coordinate radius-one bounds define a compact subset of the actual
relative spectrum by its closed product evaluation embedding and the
prescribed-base closed locus. Composing the genuine inverse affine chart
homeomorphism with its support-open subtype inclusion is continuous.
Its image of that real compact subset is therefore compact without any
separation assumption on the ambient space. When the true center belongs
to $U$, its original support also belongs to $U$, and the actual proper
lift proves all genuine section values at most one. The entire affine
chart represents the point with those original section values, giving
actual bounded-spectrum membership and recovering the whole native point
under the inverse chart.
\end{proof}

\begin{theorem}[Actual finite bounded cover and proper compactness]
\label{scheme:proper-native-compactness}
\lean{EnvelopeProperty.Berkovich.RelativeSchemeValuedPoint.iUnion_boundedAffineImage_eq_univ,
EnvelopeProperty.Berkovich.RelativeSchemeValuedPoint.isCompact_univ_of_finite_affine_cover,
EnvelopeProperty.Berkovich.RelativeSchemeValuedPoint.compactSpace_of_isProper}
\leanok
\uses{scheme:compact-bounded-native-affine-images,scheme:proper-native-center}
A genuine affine scheme cover covers the entire native prescribed-base
point space by the preceding compact bounded affine images. Any actual
finite such cover proves compactness. Genuine quasi-compactness of the
proper scheme supplies a finite affine subcover, so its entire native
prescribed-base point space is compact.
\end{theorem}
\begin{proof}
For any native point, its actual model center lies in a member of the
scheme cover. The first lemma puts that whole point in the corresponding
compact bounded affine image. This proves equality of the actual image
union with the whole native space. A finite union of real compact images
is compact in any topological space. Finally the affine base scheme
$\operatorname{Spec}R$ is quasi-compact, and the actual proper structure
map is quasi-compact, so the source scheme is quasi-compact. Mathlib's
genuine affine-basis finite-cover theorem supplies the required finite
family; applying the preceding union argument proves the final compact-space
statement. The finite affine cover is produced from the actual scheme,
not supplied as a certificate encoding native compactness.
\end{proof}
